\section{State of the field}

VOI is established in statistical decision analysis and health-economic
evaluation~\cite{claxton1999irrelevance,briggs2006decision}. R tooling
includes the Bayesian cost-effectiveness analysis (BCEA)
ecosystem~\cite{baio2017bcea}, while methodological work has made EVPPI and EVSI more
practical using regression estimators~\cite{strong2014evppi,strong2015evsi}.
Methods also exist for structural uncertainty, heterogeneity, distributional
cost-effectiveness, implementation, and adaptive trial
design~\cite{jackson2011structural,espinoza2014heterogeneity,asaria2016distributional,fenwick2008implementation,flight2022sequential}.
These are complementary advances. Reimplementing their names does not by
itself constitute a methodological contribution, and \texttt{voiage} does not
claim statistical superiority over specialist implementations.

These concerns can be evaluated in separate analyses even when they act on the
same decision. Research may have little
value if implementation is weak; subgroup optimization can conflict with an
equity-weighted objective; evidence can be precise in a source population but
poorly transportable; and waiting can preserve option value when commitment is
costly to reverse. Treating each as unrelated post-processing makes it hard to
compare their assumptions or combine them in an evidence plan.

The \texttt{voiage} implementation uses a shared, maturity-labelled problem
architecture for these value questions. Stable
classical estimands and fixture-backed frontier methods use common strategy,
value, diagnostic, and reporting semantics. The shared representation is
intended to support questions about whether the limiting uncertainty is a
parameter, study, structure, subgroup, perspective, implementation process,
data source, or decision stage without changing the basic representation of
the decision. This is an architectural design choice, not evidence of
methodological superiority. Several frontier formulations still require
empirical validation, methodological comparison, and cross-language parity
before stable promotion.
